\documentclass[tikz,border=3pt]{standalone}
\usepackage[T1]{fontenc}
\usepackage{lmodern}
\usepackage{pgfplots}
\usepackage{pgfplotstable}

\pgfplotsset{compat=1.18}

% 在仓库根目录运行：make -C paper figure-training-targets
% 图中所有柱高均从下面的 CSV 文件读取，不需要手工同步两份数据。
\newcommand{\ordinaldatafile}{../../analysis/derived/training_target_ordinal_means.csv}
\newcommand{\binarydatafile}{../../analysis/derived/training_target_binary_means.csv}

% 四种训练目标 + 推理接口组合的固定配色。
\definecolor{sods}{HTML}{7A2430}
\definecolor{rsds}{HTML}{B45D68}
\definecolor{sorf}{HTML}{174D68}
\definecolor{rsrf}{HTML}{5E8497}

% 公共坐标轴样式。bar width 与相邻 bar shift 的差值相同，因此组内柱子紧贴。
\pgfplotsset{
  common axis/.style={
    ybar,
    axis lines*=left,
    axis y discontinuity=parallel,
    grid=major,
    xmajorgrids=false,
    grid style={draw=black!18, densely dashed, line width=0.35pt},
    axis line style={draw=black!80, line width=0.65pt},
    tick style={draw=black!80, line width=0.55pt},
    tick label style={font=\sffamily\scriptsize},
    label style={font=\sffamily\small},
    title style={font=\sffamily\small\bfseries, yshift=1pt},
    xticklabel style={font=\sffamily\scriptsize, align=center},
    enlarge x limits=0.11,
    clip=true,
  },
  SO-DS/.style={
    fill=sods,
    draw=black!80,
    line width=0.45pt,
    bar width=9pt,
    bar shift=-13.5pt,
  },
  RS-DS/.style={
    fill=rsds,
    draw=black!80,
    line width=0.45pt,
    bar width=9pt,
    bar shift=-4.5pt,
  },
  SO-RF/.style={
    fill=sorf,
    draw=black!80,
    line width=0.45pt,
    bar width=9pt,
    bar shift=4.5pt,
  },
  RS-RF/.style={
    fill=rsrf,
    draw=black!80,
    line width=0.45pt,
    bar width=9pt,
    bar shift=13.5pt,
  },
}

% 两个面板分别读取独立 CSV，因此不依赖版本相关的行过滤行为。

\begin{document}
\begin{tikzpicture}

% 左侧面板：四个序数任务。修改 width/height 可调整面板尺寸。
\begin{axis}[
  common axis,
  name=ordinal,
  width=10.3cm,
  height=6.6cm,
  title={Ordinal tasks},
  ylabel={QWK},
  ymin=0.62,
  ymax=0.82,
  ytick={0.64,0.68,0.72,0.76,0.80},
  xmin=-0.55,
  xmax=3.55,
  xtick={0,1,2,3},
  xticklabels={
    Actionability,
    {Grounding\\Specificity},
    Helpfulness,
    Verifiability
  },
]
\addplot+[SO-DS] table[
  col sep=comma,
  x=x,
  y=SO_DS,
] {\ordinaldatafile};
\addplot+[RS-DS] table[
  col sep=comma,
  x=x,
  y=RS_DS,
] {\ordinaldatafile};
\addplot+[SO-RF] table[
  col sep=comma,
  x=x,
  y=SO_RF,
] {\ordinaldatafile};
\addplot+[RS-RF] table[
  col sep=comma,
  x=x,
  y=RS_RF,
] {\ordinaldatafile};

% 无边框、无背景的四色说明，纵向位置严格位于 QWK 0.76--0.80。
\path[fill=sods, draw=black!80, line width=0.35pt]
  (axis cs:1.47,0.792) rectangle (axis cs:1.60,0.798);
\node[anchor=west, font=\sffamily\tiny] at (axis cs:1.67,0.795)
  {Score-only + Direct scoring (SO--DS)};

\path[fill=rsds, draw=black!80, line width=0.35pt]
  (axis cs:1.47,0.782) rectangle (axis cs:1.60,0.788);
\node[anchor=west, font=\sffamily\tiny] at (axis cs:1.67,0.785)
  {Rationale-supervised + Direct scoring (RS--DS)};

\path[fill=sorf, draw=black!80, line width=0.35pt]
  (axis cs:1.47,0.772) rectangle (axis cs:1.60,0.778);
\node[anchor=west, font=\sffamily\tiny] at (axis cs:1.67,0.775)
  {Score-only + Rationale-first (SO--RF)};

\path[fill=rsrf, draw=black!80, line width=0.35pt]
  (axis cs:1.47,0.762) rectangle (axis cs:1.60,0.768);
\node[anchor=west, font=\sffamily\tiny] at (axis cs:1.67,0.765)
  {Rationale-supervised + Rationale-first (RS--RF)};
\end{axis}

% 右侧面板：三个二分类任务。xshift 控制两个面板之间的距离。
\begin{axis}[
  common axis,
  at={(ordinal.east)},
  anchor=west,
  xshift=1.05cm,
  width=8.5cm,
  height=6.6cm,
  title={Binary tasks},
  ylabel={Macro-F1},
  ymin=0.60,
  ymax=1.03,
  ytick={0.60,0.70,0.80,0.90,1.00},
  xmin=-0.55,
  xmax=2.55,
  xtick={0,1,2},
  xticklabels={
    Coherence,
    {Positioning\\Check},
    {Positioning\\Type}
  },
]
\addplot+[SO-DS] table[
  col sep=comma,
  x=x,
  y=SO_DS,
] {\binarydatafile};
\addplot+[RS-DS] table[
  col sep=comma,
  x=x,
  y=RS_DS,
] {\binarydatafile};
\addplot+[SO-RF] table[
  col sep=comma,
  x=x,
  y=SO_RF,
] {\binarydatafile};
\addplot+[RS-RF] table[
  col sep=comma,
  x=x,
  y=RS_RF,
] {\binarydatafile};
\end{axis}

\end{tikzpicture}
\end{document}
